%Text automatically generated by txt2latex
\section{DESCRIPTION}
codata is a Fortran library providing the fundamental physical
constants according to CODATA. 
A C API allows usage from C, or can be used as a basis for other wrappers.
A Python wrapper allows easy usage from Python.

\section{FEATURES}
It provides:
\begin{itemize}
\item codata constants 2022
\item codata constants 2018
\item codata constants 2014
\item codata constants 2010
\end{itemize}

\section{INSTALLATION}
See INSTALL.

\section{DOCUMENTATION}
See docs/ for detailed documentation or directly online at 
https://milanskocic.github.io/codata

\section{NOTES FOR FORTRAN USERS}
The codata constant from 2022 were integrated in the Fortran stdlib.
See https://github.com/fortran-lang/stdlib/releases/tag/v0.7.0.
If you need older values for the CODATA constants you can use this library 
otherwise use the standard library.
To use codata within your fpm(1) (https://github.com/fortran-lang/fpm)
project, add the following lines to your file:

\begin{verbatim}
     [dependencies]
     codata = { git="https://github.com/MilanSkocic/codata.git" }

\end{verbatim}
If you only need sources for the codata constants that you can integrate
directly in your sources you may be interested by this repo
https://github.com/vmagnin/fundamental\_constants.

\section{LICENSE}
MIT
